% MTH321 Project 1, Direction 4 (Financial SDE) -- Team F talk.
% Build (from slides/): latexmk -pdf TeamF_Slides.tex    Figures: python make_slide_figures.py (reads ../results/*.csv)
% Numbers: generated/macros.tex (written by ../code/report_tables.py; same macros as the report) and the report tables / CSVs.
\documentclass[aspectratio=169,10pt,t]{beamer}
% =====================================================================
% polish_preamble.tex -- Team F talk preamble: palette, hairlines, theorem style, sizes
% (body 10pt, frame titles 12pt, tables 9pt, captions 8pt), old-style figures in text,
% lining figures in math, quiet section navigation in the footer.
% =====================================================================
\usepackage[T1]{fontenc}
\usefonttheme{serif}
\usefonttheme{professionalfonts}
\usepackage[sc,osf]{mathpazo}       % Palatino; old-style figures in text, lining figures in math
\usepackage{amsmath,amssymb,array}
\linespread{1.06}

% ---------- palette --------------------------------------------------------
\definecolor{ink}{HTML}{1F2933}
\definecolor{accent}{HTML}{23426B}
\definecolor{muted}{HTML}{4B5563}   % secondary text, lighter than ink
\definecolor{hairline}{HTML}{C9CDD2}
\definecolor{note}{HTML}{374151}     % page-bottom notes and captions (darker than muted)
\definecolor{navcur}{HTML}{C2642B}   % current section in the footer nav (muted burnt orange)
\definecolor{hlY}{HTML}{FFEB99}
\definecolor{hlG}{HTML}{CDEBC5}
\definecolor{hlR}{HTML}{F9D3C4}
\definecolor{boxred}{HTML}{B03A2E}
\definecolor{lineblue}{HTML}{2F5D8C}
\setbeamercolor{normal text}{fg=ink,bg=white}
\setbeamercolor{background canvas}{bg=white}
\setbeamercolor{structure}{fg=accent}
\setbeamercolor{alerted text}{fg=boxred}
\setbeamercolor{frametitle}{fg=accent,bg=}
\setbeamercolor{footline}{fg=accent}   % footer in the structure navy
\setbeamercolor{itemize item}{fg=muted}
\setbeamercolor{enumerate item}{fg=muted}

% ---------- sizes -------------------------------------------------------------
\setbeamerfont{frametitle}{size=\fontsize{12}{14.5}\selectfont}
\setbeamerfont{footline}{size=\fontsize{7.5}{9}\selectfont}

% ---------- layout --------------------------------------------------------------
\setbeamersize{text margin left=8mm,text margin right=8mm}
\setbeamertemplate{navigation symbols}{}
\setbeamertemplate{headline}{}
\setbeamertemplate{frametitle}{%
  \vskip 4.2mm
  {\usebeamerfont{frametitle}\usebeamercolor[fg]{frametitle}\insertframetitle\par}
  \vskip 1.3mm{\color{hairline}\rule{\textwidth}{0.4pt}}\vskip -2.6mm}
% quiet agenda in the footer: the current part in accent, the others in muted small caps
\newcommand{\navitem}[2]{\ifnum\value{section}=#1{\color{navcur}\textsc{#2}}\else\textsc{#2}\fi}   % current part in burnt orange
\newcommand{\navsep}{\hspace{0.9em}}
\setbeamertemplate{footline}{%
  \usebeamerfont{footline}\usebeamercolor[fg]{footline}%
  \hbox to \paperwidth{\hskip 8mm\rlap{\textsc{Financial SDEs}\ \textperiodcentered\ \textsc{Team F}}\hfill
    \navitem{1}{setting}\navsep\navitem{2}{gbm}\navsep\navitem{3}{non-affine}\navsep
    \navitem{4}{checks}\navsep\navitem{5}{close}\hfill
    \llap{\insertframenumber\,/\,\inserttotalframenumber}\hskip 8mm}\vskip 3.6mm}
\addtobeamertemplate{frame begin}{}{\raggedright\hyphenpenalty=10000\exhyphenpenalty=10000}
\hyphenpenalty=10000 \exhyphenpenalty=10000
\setbeamertemplate{itemize item}{\raisebox{0.12ex}{\textendash}}
\setbeamertemplate{enumerate item}{\insertenumlabel.}
\setlength{\leftmargini}{1.4em}

% ---------- theorem-like: small-caps label, no coloured blocks ------------------
\setbeamertemplate{theorems}[normal font]
\setbeamertemplate{theorem begin}{\par\vskip 1.0ex\noindent
  {\upshape\scshape\color{accent}\inserttheoremname}%
  \ifx\inserttheoremaddition\empty\else\ {\color{muted}(\inserttheoremaddition)}\fi.\ \itshape}
\setbeamertemplate{theorem end}{\par\vskip 0.8ex}
\newtheorem*{proposition*}{Proposition}
\newtheorem*{remark*}{Remark}

% ---------- no emphasis marks: no boxes, underlines or highlights -------------------
\newcommand{\dfn}[1]{{\color{accent}#1}}
\newcommand{\notesize}{\fontsize{8.5}{10.6}\selectfont}         % notes/captions 8.5pt
\newcommand{\notesizesmall}{\fontsize{7.5}{9.4}\selectfont}    % AI-use block 7.5pt
\newcommand{\aside}[1]{{\color{note}\notesize\upshape #1}}
\newcommand{\figcap}[1]{{\color{note}\notesize\upshape #1\par}}   % \par inside: footnotesize leading

% ---------- tables and display spacing ---------------------------------------------
\usepackage{booktabs}
\setlength{\aboverulesep}{0.3ex}\setlength{\belowrulesep}{0.45ex}
\setlength{\heavyrulewidth}{0.6pt}\setlength{\lightrulewidth}{0.35pt}\setlength{\cmidrulewidth}{0.35pt}
\renewcommand{\arraystretch}{1.1}
\newcommand{\tabfont}{\small\fontfamily{pplx}\selectfont}   % 9pt, lining figures so table columns align
\newcommand{\stack}[1]{\begin{tabular}[t]{@{}l@{}}#1\end{tabular}}   % manual line break inside a cell
\makeatletter
\g@addto@macro\normalsize{\setlength\abovedisplayskip{6pt}\setlength\belowdisplayskip{6pt}%
  \setlength\abovedisplayshortskip{4pt}\setlength\belowdisplayshortskip{4pt}}
\makeatother

% ---------- layout rhythm ---------------------------------------------------
% every frame body starts \bodytop below the title rule; footnote-style asides sit \asidegap above the footer
\newlength{\bodyskip}\setlength{\bodyskip}{1.6mm}
\newcommand{\bodytop}{\vspace*{\bodyskip}}
\newcommand{\tabletop}{\vspace*{\bodyskip}\vspace*{2.5mm}}   % table top rule level with the cap height of text pages
\newcommand{\closing}[1]{\par\vspace{0pt plus 1filll}{\color{note}\notesize\upshape #1\par}\vspace{2.6mm}}   % filll: beats beamer's [t] bottom fill, so every aside sits on the same line
\newcommand{\blockgap}{\par\vspace{0.9em}}     % between body blocks (table -> text, list -> text)
\newcommand{\colgap}{\par\vspace{0.75em}}      % between paragraphs inside the text column of figure frames
\newcommand{\lead}[1]{{\upshape\scshape\color{accent}#1}}   % small-caps run-in label (Goal, Remark)

% ---------- stroke weight: text drawn fill+stroke with a hairline stroke in the text colour ----
% Palatino stays Palatino; glyphs gain ~0.1pt per side, so the deck reads better when projected.
% Text remains real text (selectable, searchable); colours unchanged (xcolor sets fill and stroke).
\usepackage{pdfrender}
\newcommand{\strokeweight}{0.15pt}
\AtBeginDocument{\pdfrender{TextRenderingMode=FillStroke,LineWidth=\strokeweight,LineJoinStyle=1}}

\graphicspath{{figures/}}
\newcommand{\E}{\mathbb{E}}
% figure frame: figure (natural size, as drawn) and caption left, text right; tops aligned
\newcommand{\figframe}[3]{%
  \begin{columns}[T,onlytextwidth]
  \column{0.49\textwidth}
    \includegraphics[height=\figH]{#1}\par\vspace{1.2mm}
    {\raggedright\figcap{#2}\par}
  \column{0.05\textwidth}
  \column{0.46\textwidth}\raggedright\vspace{\figtextdrop}
    #3
  \end{columns}}
\newlength{\figH}\setlength{\figH}{47mm}% all four figures at one height (same scale, same font size)
\newlength{\figtextdrop}\setlength{\figtextdrop}{0.3mm}
% series colours = figure colours (slide_rcparams.py: ACCENT #23426B navy, RED #B03A2E brick); label cells only
\definecolor{figNavy}{HTML}{23426B}\definecolor{figRed}{HTML}{B03A2E}
\newcommand{\serA}[1]{{\color{figNavy}#1}}\newcommand{\serB}[1]{{\color{figRed}#1}}
\newcommand{\slopetab}[1]{\begin{tabular}{@{}l@{\hspace{1.1em}}l@{}}#1\end{tabular}}
\newcommand{\idfig}{\fontfamily{pplx}\itshape}   % student IDs: Palatino italic, lining (equal-height, tabular) figures
\newcommand{\emLocalOrderFine}{0.500}
\newcommand{\milLocalOrderFine}{0.999}
\newcommand{\emLocalOrderCoarse}{0.441}
\newcommand{\milLocalOrderCoarse}{0.916}
\newcommand{\analyticBiasHalf}{-0.0646}
\newcommand{\analyticInCI}{18}
\newcommand{\eTwoSlopeA}{1.021}
\newcommand{\eTwoSlopeSEA}{0.069}
\newcommand{\eTwoSlopeB}{1.013}
\newcommand{\eTwoSlopeSEB}{0.083}
\newcommand{\eOneUncoupled}{23.62}
\newcommand{\eOneUncSlopeEM}{-0.0004}
\newcommand{\eOneUncSlopeMil}{-0.0007}
\newcommand{\eOneCoupledSlopeEM}{0.4941}
\newcommand{\eOneCoupledSlopeMil}{0.9906}
\newcommand{\eTwoPayoffTwo}{0.0724}
\newcommand{\eTwoPayoffTwoSE}{0.0046}
\newcommand{\eTwoMeanTwo}{0.0024}
\newcommand{\eTwoMeanTwoSE}{0.0053}
\newcommand{\eTwoOracleAll}{8}
\newcommand{\eTwoOracleN}{8}
\newcommand{\eTwoPairedAll}{7}
\newcommand{\eTwoPairedN}{7}
\newcommand{\eTwoCommonRange}{2,4,8,16}
\newcommand{\eTwoCommonSlopeA}{1.021}
\newcommand{\eTwoCommonSlopeSEA}{0.069}
\newcommand{\eTwoCommonSlopeB}{1.001}
\newcommand{\eTwoCommonSlopeSEB}{0.067}
\newcommand{\eTwoGateThirtyTwo}{0.499}
\newcommand{\eTwoRefGap}{0.00044}
\newcommand{\eTwoRefGapSE}{0.00020}
\newcommand{\eTwoMeanFiveTwelve}{0.00040}
\newcommand{\eTwoMeanFiveTwelveSE}{0.00024}
\newcommand{\eTwoMinPairedSE}{0.00024}
\newcommand{\eTwoSmallestFit}{0.0045}
\newcommand{\eTwoSmallestFitN}{32}
\newcommand{\eTwoCallRef}{14.574}
\newcommand{\eTwoCallRefSE}{0.034}
\newcommand{\anaVarEM}{450.809}
\newcommand{\anaVarMilstein}{450.845}
\newcommand{\anaMean}{105.127}
\newcommand{\emMean}{105.12035}
\newcommand{\exactMean}{105.12711}
\newcommand{\exactVar}{451.02881}
\newcommand{\quantInside}{10}
\newcommand{\quantTotal}{10}
\newcommand{\payoffZeroIn}{10}
\newcommand{\payoffTotal}{10}
\newcommand{\refGap}{0.04022}
\newcommand{\refGapSE}{0.00015}
\newcommand{\refFinest}{0.15495}
\newcommand{\refRatio}{25.96}
\newcommand{\refChange}{0.00552}
\newcommand{\refSlopeChange}{0.01153}
\newcommand{\costEMN}{512}
\newcommand{\costMilsteinN}{8}
\newcommand{\wcCEM}{34.35}
\newcommand{\wcCMil}{0.336}
\newcommand{\wcCRatio}{102}
\newcommand{\wcCSlowEM}{6.6}
\newcommand{\wcCSlowMil}{1.7}
\newcommand{\wcEM}{5.18}
\newcommand{\wcMil}{0.199}
\newcommand{\wcRatio}{26}
\newcommand{\wcPaths}{10{,}000}
\newcommand{\wcRepeats}{7}
\newcommand{\wcEMEight}{0.066}
\newcommand{\wcPerStepEMEight}{8.22\times10^{-10}}
\newcommand{\wcEMFiveTwelve}{5.179}
\newcommand{\wcPerStepEMFiveTwelve}{1.01\times10^{-9}}
\newcommand{\wcMilEight}{0.199}
\newcommand{\wcPerStepMilEight}{2.49\times10^{-9}}
\newcommand{\wcMilFiveTwelve}{13.232}
\newcommand{\wcPerStepMilFiveTwelve}{2.58\times10^{-9}}
\newcommand{\wcStepRatioEight}{3.02}
\newcommand{\wcStepRatioFiveTwelve}{2.55}
\newcommand{\wcPython}{3.13.5}
\newcommand{\wcNumpy}{2.5.3}
\newcommand{\wcMachine}{x86\_64}
\newcommand{\eOneBGbmEMOkSlope}{0.501}
\newcommand{\eOneBGbmEMBadSlope}{0.501}
\newcommand{\eOneBGbmMilOkSlope}{0.993}
\newcommand{\eOneBGbmMilBadSlope}{0.993}
\newcommand{\eOneBNonaffOkSlope}{0.519}
\newcommand{\eOneBNonaffBadSlope}{0.028}
\newcommand{\eOneBNonaffOkErr}{0.150}
\newcommand{\eOneBNonaffBadErr}{1.839}
\newcommand{\eOneBNonaffBadErrCoarse}{1.985}
\newcommand{\eOneBPaths}{10{,}000}

\title{Numerical solution and validation of financial SDEs}
\hypersetup{pdftitle={Numerical solution and validation of financial SDEs},
  pdfauthor={Team F: Kunze Li, Che Li, Zekun Wang, Yixuan Liu, Yawen Wu, Ruitong Huang},
  pdfsubject={MTH321 Project 1, Direction 4: Financial SDE},
  pdfkeywords={Euler-Maruyama, Milstein, geometric Brownian motion, non-affine stochastic volatility, log-Euler}}

\begin{document}
\section{Setting}
% =========================================================== 1  cover
% Cover page number: the frame is [plain] (no footline, so the title block keeps its
% vertical balance); "1 / 14" is drawn on the background layer with exactly the
% footline's font, colour and geometry (right margin 8mm, 3.6mm above the bottom edge).
{\setbeamertemplate{background}{\vbox to \paperheight{\vfill
  \usebeamerfont{footline}\usebeamercolor[fg]{footline}%
  \hbox to \paperwidth{\hfill\llap{\insertframenumber\,/\,\inserttotalframenumber}\hskip 8mm}\vskip 3.6mm}}
% XJTLU crest (shield only), vector PDF: figures/xjtlu_crest.pdf (+ .svg).
%   Source: official XJTLU logo SVG on xjtlu.edu.cn,
%     https://www.xjtlu.edu.cn/wp-content/uploads/2024/01/cn-footer-logo-light.svg
%   as used on https://www.xjtlu.edu.cn/en/about/overview/xjtlu-identity .
%   Shield paths extracted unchanged; the reversed-out (light) colours were swapped
%   to the positive navy shield (#010544) with white charges, as in the official
%   logo on that page. Converted SVG -> PDF with rsvg-convert (vector, no raster).
%   Licence: the XJTLU logo is a registered trademark of Xi'an Jiaotong-Liverpool
%   University (see https://www.xjtlu.edu.cn/en/about/brand-resources); used here
%   only to identify the institution on a non-commercial student coursework talk.
% Placement: height 14mm, right edge on the 8mm right margin, top aligned with the
% small-cap height of the cover's top text line.
\newlength{\crestcap}
\begin{frame}[plain]
\vspace*{0pt plus 1filll}% block centred optically (top : bottom = 1 : 1.25)
\hspace*{6mm}\begin{minipage}{0.86\textwidth}\raggedright
{\color{note}\notesize\settoheight{\crestcap}{\textsc{mth}}%
\leavevmode\rlap{\hspace*{\dimexpr\paperwidth-22mm\relax}\llap{\raisebox{\dimexpr\crestcap-14mm\relax}[0pt][0pt]{\includegraphics[height=14mm]{xjtlu_crest}}}}%
\textsc{mth321}\ \textperiodcentered\ \textsc{project 1}\ \textperiodcentered\ \textsc{direction 4: financial sde}}\par\vspace{5mm}
{\color{accent}\fontsize{13}{16.5}\selectfont Numerical solution and validation of financial SDEs\par}\vspace{1.6mm}
{\itshape Geometric Brownian motion and a non-affine stochastic-volatility model\par}\vspace{3.5mm}
{\color{hairline}\rule{0.42\textwidth}{0.4pt}}\par\vspace{3.2mm}
{\textsc{Team F}}\par\vspace{1.6mm}
{\small
\begin{tabular}{@{}l@{\hspace{1.6em}}>{\idfig}l@{\hspace{1.6em}}>{\color{note}\itshape}l@{}}
{\color{note}\footnotesize\textsc{member}} & {\normalfont\color{note}\footnotesize\textsc{student id}} & {\upshape\footnotesize\textsc{role}}\\[0.25em]
Kunze Li & 2363973 & Project Manager\\
Che Li & 2362072 & Algorithm Implementation\\
Zekun Wang & 2363718 & Validation \& Visualization\\
Yixuan Liu & 2361140 & Testing \& Validation\\
Yawen Wu & 2360749 & Mathematical Theory\\
Ruitong Huang & 2360294 & Visualization \& Report\\
\end{tabular}}\par\vspace{4mm}
{\color{accent}\footnotesize Xi'an Jiaotong-Liverpool University \textperiodcentered\ October 2026}\par
\end{minipage}
\vspace*{0pt plus 1.25filll}
\end{frame}}

% =========================================================== 2  questions
\begin{frame}{Two questions, answered against two kinds of reference}
\bodytop
\begin{enumerate}\setlength{\itemsep}{0.8em}
\item For geometric Brownian motion, whose exact solution is known, do coupled Monte~Carlo experiments recover the
  theoretical strong and weak orders of Euler--Maruyama and Milstein once sampling uncertainty is quantified?
\item For a non-affine stochastic-volatility model without an exact transition law, does Euler--Maruyama refine at half order
  against nested numerical references, and can its weak payoff bias be resolved statistically?
\end{enumerate}
\vspace{0.95em}
{\leftskip=\leftmargini\lead{Why both.}\ Strong error is pathwise, weak error concerns expectations; their rates differ.\par}% content: report Sec. 1.1
\vspace{0.6em}
\hspace*{\leftmargini}\lead{Goal.}\ A trend counts only where it exceeds the Monte Carlo uncertainty.
\closing{In Direction 4, stability regions and adaptive steps are replaced by the exact--EM--Milstein comparison, coupled strong and
weak convergence, Monte Carlo confidence intervals and nested-grid validation; we add an analytic positivity analysis.}
\end{frame}

% =========================================================== 3  GBM schemes
\begin{frame}{GBM: the exact solution and two one-step schemes}
\bodytop\vspace{1.0mm}\vspace{-\abovedisplayskip}\vspace{-\baselineskip}
\begin{align*}
dS_t&=\mu S_t\,dt+\sigma S_t\,dW_t, & S_T&=S_0\exp\bigl[(\mu-\tfrac12\sigma^2)T+\sigma W_T\bigr],\\[0.2em]
S^{\mathrm{EM}}_{n+1}&=S_n\bigl(1+\mu h+\sigma\Delta W_n\bigr), &
S^{\mathrm{Mil}}_{n+1}&=S_n\bigl[1+\mu h+\sigma\Delta W_n+\tfrac12\sigma^2\bigl((\Delta W_n)^2-h\bigr)\bigr].
\end{align*}
Errors are measured on the same Brownian path,
\[
e_s(h)=\E\bigl|S^{(h)}_N-S_T\bigr|,\qquad e_w(h)=\bigl|\E S^{(h)}_N-\E S_T\bigr|,
\]
and theory predicts strong order $\tfrac12$ for Euler--Maruyama, $1$ for Milstein, and weak order $1$ for both.
\par\vspace{0.75em}% content: report Sec. 2.7 (positivity replaces stability)
\lead{Positivity.}\ The Euler--Maruyama factor $1+\mu h+\sigma\sqrt h\,Z$ can be negative for every $h$, the Milstein
factor only if $(\sigma^2-2\mu)h>1$; here $\sigma^2-2\mu=-0.06$ and $p_{\mathrm{EM}}(1/2)\approx2.12\times10^{-13}$.
\closing{$S_0=100$, $\mu=0.05$, $\sigma=0.20$, $T=1$; $M=2\times10^6$ paths, seed $32104$; $N=2,4,\dots,512$; coarse
increments are block sums of one $512$-step master path.}
\end{frame}

% =========================================================== 4  non-affine model
\begin{frame}{A non-affine volatility model with non-commutative noise}
\bodytop\vspace{1.8mm}\vspace{-\abovedisplayskip}\vspace{-\baselineskip}
\begin{gather*}
dX_t=\bigl(\mu-\tfrac12g(Y_t)^2\bigr)\,dt+g(Y_t)\,dW^{(1)}_t,\qquad
dY_t=\kappa(\theta-Y_t)\,dt+\xi\sqrt{1+Y_t^2}\,dW^{(2)}_t,\\[0.2em]
g(y)=\sigma_{\min}+\frac{\sigma_{\max}-\sigma_{\min}}{1+e^{-y}},\qquad d\langle W^{(1)},W^{(2)}\rangle_t=\rho\,dt,\qquad S_t=e^{X_t}.
\end{gather*}
\begin{proposition*}[non-commutativity]
With $q(y)=\xi\sqrt{1+y^2}$, the diffusion columns $b_1,b_2$ of $(X,Y)$ for independent drivers $B^{(1)}=W^{(1)}$, $B^{(2)}$ satisfy
\[
L^2b_1-L^1b_2=\bigl(\sqrt{1-\rho^2}\,q(Y)\,g'(Y),\,0\bigr)^{\!\top}\neq0.
\]
\end{proposition*}
Strong order one would in general need L\'evy areas (commutator $X$-component $0.0433$ at $Y=\theta$); we use
log-Euler--Maruyama, for which $S_{n+1}=e^{X_{n+1}}>0$.
\closing{Model: $\mu=0.05$, $\kappa=2$, $\theta=-0.2$, $\xi=0.6$, $\sigma_{\min}=0.10$, $\sigma_{\max}=0.50$.\\
Initial values, horizon and correlation: $S_0=100$, $Y_0=0$, $T=1$, $\rho=-0.7$; payoff $(S_T-100)^+$.}
\end{frame}

% =========================================================== 5  oracles
\begin{frame}{Two exact identities serve as oracles for the code}
\bodytop\vspace{-0.6mm}
\begin{proposition*}[GBM mean bias]
For Euler--Maruyama and Milstein, $\E S_N=S_0(1+\mu h)^N$, hence
\[
\delta(h)=S_0\bigl[(1+\mu h)^{T/h}-e^{\mu T}\bigr]=-\tfrac12S_0e^{\mu T}\mu^2T\,h+O(h^2).
\]
\end{proposition*}
For our parameters $\delta(1/2)=\analyticBiasHalf$; with $\mu=0$ the bias would vanish and $\varphi(s)=s$ could not reveal a
weak rate.% content: report Sec. 2.5
\par\vspace{0.85em}
\begin{proposition*}[log-EM preserves the mean]
For the non-affine log-Euler--Maruyama scheme and every $h=T/N$, \ $\E S_N=S_0e^{\mu T}=\E S_T$.
\end{proposition*}
\vspace{1.4em}
\noindent\lead{Remark.}\ The identity payoff therefore carries no discretisation bias in the non-affine model; its weak error
must be measured with a nonlinear payoff, here the call $(S_T-100)^+$.
\closing{Second identity as a code oracle: the paired difference $S^{(h)}_N-S^{(h_{\mathrm{ref}})}_N$ must be zero within error.}% content: report Remark 2.5, Sec. 4.6 (oracle passed)
\end{frame}

\section{GBM}
% =========================================================== 6  GBM strong
\begin{frame}{On GBM the strong orders one half and one are reproduced}
\note{Data: results/gbm\_strong.csv}
\bodytop
\figframe{gbm_strong.pdf}
{Strong error against the exact coupled path, $M=2\times10^6$. Filled points ($N\ge64$) fitted; $95\%$ bars within
markers; dashed: slopes $\tfrac12$ and~$1$.}
{Fitted slopes, $\pm$ one jackknife SE:\par\vspace{0.4em}
\slopetab{\serA{Euler--Maruyama} & $0.4991\pm0.0004$\\ \serB{Milstein} & $0.9979\pm0.0002$}\colgap
Local orders rise monotonically to $\emLocalOrderFine$ and $\milLocalOrderFine$; at $h=1/512$ Milstein is about $60\times$
more accurate.\colgap
The coarse-grid deficit is pre-asymptotic curvature, so only $N\ge64$ is fitted; the small gap to $\tfrac12$ and $1$ is
finite-grid curvature, which the SE does not measure.}% content: report Sec. 4.1
\end{frame}

% =========================================================== 7  GBM weak
\begin{frame}{The weak bias of GBM is first order, as the analytic bias predicts}
\note{Data: results/gbm\_weak.csv; moments and quantiles: results/gbm\_moments.csv, results/gbm\_distribution.csv}
\bodytop
\figframe{gbm_weak.pdf}
{Top: absolute bias and analytic $|\delta(h)|$; bottom: signed bias with $95\%$ intervals. Filled points fitted;
markers offset horizontally for visibility.}
{Weak slopes:\par\vspace{0.4em}
\slopetab{\serA{Euler--Maruyama} & $1.0100\pm0.0440$\\ & {\color{muted}\footnotesize $N=2$--$64$}\\[0.15em]
\serB{Milstein} & $0.9981\pm0.0004$\\ & {\color{muted}\footnotesize $N=2$--$512$}}\colgap
All $\analyticInCI$ signed intervals contain $\delta(h)$.\colgap
At $N=256$, moments and all ten quantile intervals agree with the exact law within Monte Carlo error.}
\end{frame}

\section{Non-affine}
% =========================================================== 8  non-affine strong
\begin{frame}{Log-EM shows an empirical half order against two references}
\note{Data: results/nonaffine\_strong.csv}
\bodytop
\figframe{nonaffine_strong.pdf}
{Coupled discrepancy to Euler--Maruyama references, $M=5\times10^4$. All points fitted, offset
horizontally; $95\%$ bars within markers; dashed: slope~$\tfrac12$.}
{Slopes on $N=32$--$512$:\par\vspace{0.4em}
\slopetab{\serA{ref.\ $4096$} & $0.5205\pm0.0019$\\ \serB{ref.\ $8192$} & $0.5090\pm0.0018$}\colgap
Local orders rise at the finest grids (to $0.558$ against $4096$, $0.530$ against $8192$): a reference only $8\times$
finer shares part of the coarse error, so the finer reference is closer to~$\tfrac12$.\colgap% content: report Sec. 4.4, Table 6
The references differ by $\refGap\pm\refGapSE$ ($\refRatio\%$ of the value at $N=512$): these are discrepancies, not exact
errors.}
\end{frame}

% =========================================================== 9  non-affine weak, main run
\begin{frame}{With $5\times10^4$ paths the non-affine payoff bias is not resolved}
\note{Data: results/nonaffine\_weak.csv}
\bodytop
\figframe{nonaffine_weak.pdf}
{Paired call-payoff differences against both references; top: absolute, bottom: signed with $95\%$ intervals.
All points open (none fitted), offset horizontally.}
{All ten intervals contain zero (paired, $95\%$; SE about $0.003$ at $N=32$ and $0.0008$ at $N=512$).\colgap
No weak slope is reported: no grid passes the resolution rule in this run.\colgap
Common random numbers make the paired SE small, but for $N\ge32$ the bias is smaller still; a slope fitted here would
measure sampling error, not the scheme. E2 resolves the bias on coarser grids.}% content: report Sec. 4.5, Fig. 5
\end{frame}

% E2 slopes on p10 written to 4 dp as in report Table tab:slopes (macros give 3 dp), so p10 and p14 agree
% =========================================================== 10  E2
\begin{frame}{With $4\times10^5$ paths on coarse grids the payoff bias is resolved}
\note{Data: results/extensions/ext\_nonaffine\_coarse.csv}
\tabletop
{\tabfont
\begin{tabular*}{\textwidth}{@{}r@{\extracolsep{\fill}}lll@{}}
\toprule
$N$ & \stack{payoff difference,\\ ref.\ $512$} & \stack{payoff difference,\\ ref.\ $1024$} & \stack{mean difference,\\ ref.\ $1024$}\\
\midrule
2  & $0.0719\pm0.0046$ & $0.0724\pm0.0046$ & $\phantom{-}0.0024\pm0.0053$\\
4  & $0.0402\pm0.0033$ & $0.0406\pm0.0033$ & $\phantom{-}0.0014\pm0.0038$\\
8  & $0.0174\pm0.0023$ & $0.0178\pm0.0023$ & $-0.0009\pm0.0027$\\
16 & $0.0090\pm0.0016$ & $0.0094\pm0.0016$ & $\phantom{-}0.0005\pm0.0019$\\
32 & $0.0040\pm0.0011^{\circ}$ & $0.0045\pm0.0011$ & $\phantom{-}0.0007\pm0.0013$\\
64 & $0.00196\pm0.00077^{\circ}$ & $0.00241\pm0.00079^{\circ}$ & $\phantom{-}0.00071\pm0.00092$\\
\bottomrule
\end{tabular*}\par}
\blockgap
The weak slopes are $1.0209\pm0.0685$ (ref.\ $512$, $N=2$--$16$) and $1.0127\pm0.0829$
\mbox{(ref.\ $1024$, $N=2$--$32$)}. The mean difference is zero within error, as the identity requires.
\closing{$M=4\times10^5$; $\pm$ one SE; ${}^\circ$ not resolved by the fit rule; grid set fixed after a first run with $N=1$.}
\end{frame}

\section{Checks}
% =========================================================== 11  E1 / E1b
\begin{frame}{The strong rates disappear when the coupling is removed}
\note{Data: results/extensions/ext\_uncoupled\_strong.csv, results/extensions/ext\_misaligned.csv}
\tabletop
{\tabfont
\begin{tabular*}{\textwidth}{@{}r@{\extracolsep{\fill}}rrrr@{}}
\toprule
 & \multicolumn{2}{c}{coupled} & \multicolumn{2}{c}{uncoupled}\\
\cmidrule(r{1.4em}){2-3}\cmidrule(l{1.4em}){4-5}
$N$ & Euler--Maruyama & Milstein & Euler--Maruyama & Milstein\\
\midrule
8   & $0.8179$  & $0.10665$   & $23.573$ & $23.544$\\
32  & $0.41633$ & $0.027503$  & $23.606$ & $23.599$\\
128 & $0.20912$ & $0.006937$  & $23.615$ & $23.613$\\
512 & $0.10492$ & $0.0017351$ & $23.617$ & $23.616$\\
\addlinespace[3pt]
slope & $\eOneCoupledSlopeEM$ & $\eOneCoupledSlopeMil$ & $\eOneUncSlopeEM$ & $\eOneUncSlopeMil$\\
\bottomrule
\end{tabular*}\par}
\blockgap
Without coupling the error is the natural spread $\E|S_T-S_T'|\approx\eOneUncoupled$ (E1), $29\times$ to $225\times$ the
coupled Euler--Maruyama error. Replacing the block sum by a
strided sum leaves the GBM slopes unchanged ($\eOneBGbmEMOkSlope$, $\eOneBGbmMilOkSlope$) but lowers the non-affine slope
from $\eOneBNonaffOkSlope$ to $\eOneBNonaffBadSlope$ (E1b).
\closing{E1: GBM strong error, $M=2\times10^5$, independent seed stream. E1b: $M=10^4$.}
\end{frame}

% =========================================================== 12  cost
\begin{frame}{Cost at a matched strong-error target of $0.11$}
\note{Data: results/gbm\_strong.csv, results/extensions/ext\_wallclock.csv}
\tabletop
{\tabfont
\begin{tabular*}{\textwidth}{@{}l@{\extracolsep{\fill}}rrrrr@{}}
\toprule
Method & $N$ & strong error & upper $95\%$ & work $MN$ & solver time\\
\midrule
Euler--Maruyama & $512$ & $0.10473$ & $0.10485$ & $1{,}024{,}000{,}000$ & $\wcEM$\,ms\\
Milstein        & $8$   & $0.10658$ & $0.10672$ & $16{,}000{,}000$     & $\wcMil$\,ms\\
\bottomrule
\end{tabular*}\par}
\blockgap
Milstein reaches the target with $64\times$ fewer path updates and about $26\times$ less wall-clock time; the second ratio
is smaller because one Milstein step costs $\wcStepRatioEight$ ($N=8$) to $\wcStepRatioFiveTwelve$ ($N=512$) times an
Euler--Maruyama step.\blockgap
The target $0.11$ is about the smallest error Euler--Maruyama attains on the tested grids ($N\le512$). Random-number
generation, not counted, would also favour Milstein.% content: report Sec. 4.9
\closing{GBM, $M=2\times10^6$ for the errors. Time: solver loop only, $\wcPaths$ paths, median of $\wcRepeats$ repeats, one
machine. The $158\times$ of our first version was a row-major storage artefact, now corrected.}
\end{frame}

\section{Close}
% =========================================================== 13  limitations
\begin{frame}{Limitations}
\bodytop
\begin{itemize}\setlength{\itemsep}{0.38em}
\item The non-affine errors are discrepancies against Euler--Maruyama references, not exact~errors; the reference gap
  is $\refRatio\%$ of the finest discrepancy.
\item The E2 weak slopes rest on coarse grids; the grid set was chosen after a first run.
\item The call payoff is not smooth and $e^X$ is not globally Lipschitz, so the non-affine rates are~empirical.
\item One parameter set; positivity was analysed only analytically; no Milstein scheme with L\'evy areas.
\item Slope errors are conditional jackknife SEs; timings come from one machine.
\end{itemize}
\closing{{\color{hairline}\rule{\textwidth}{0.4pt}}\par\vspace{0.6mm}\notesizesmall
AI use (report, Appendix A). Codex: numerical library and tests, review-driven revisions (resolved weak grids, jackknife SE);
on 10 Oct.\ (GPT-6 Astra, Ultra) the weak-grid refactor, independently re-checked. ChatGPT (GPT-6, Pro: review; GPT-6.1 Sol: language
polishing). Grok Bot: LaTeX typesetting, report restructuring and code packaging; drafted the supplementary code and two
derivations.}
\end{frame}

% =========================================================== 14  summary
\begin{frame}{Summary of the observed rates}
\note{Data: report Table tab:slopes}
\tabletop
{\tabfont
\begin{tabular*}{\textwidth}{@{}l@{\extracolsep{\fill}}lcll@{}}
\toprule
Study & Method & theory & slope $\pm$ SE & fitted $N$\\
\midrule
GBM strong & Euler--Maruyama & $1/2$ & $0.4991\pm0.0004$ & 64--512\\
GBM strong & Milstein & $1$ & $0.9979\pm0.0002$ & 64--512\\
GBM weak & Euler--Maruyama & $1$ & $1.0100\pm0.0440$ & 2--64\\
GBM weak & Milstein & $1$ & $0.9981\pm0.0004$ & 2--512\\
\addlinespace[2pt]
Non-affine strong & ref.\ 4096 / 8192 & ($1/2$) & $0.5205\pm0.0019$ / $0.5090\pm0.0018$ & 32--512\\
Non-affine weak, E2 & ref.\ 512 / 1024 & ($1$) & $1.0209\pm0.0685$ / $1.0127\pm0.0829$ & 2--16 / 2--32\\
\bottomrule
\end{tabular*}\par}
\blockgap
On GBM the theory is reproduced against the exact solution; for the non-affine model the rates are empirical and depend on
numerical references. Coupling is essential.% content: report Sec. 5 (Conclusions)
\ Next steps: a third reference, multilevel Monte Carlo, and Milstein with simulated L\'evy areas.
\closing{Theory in parentheses is expected, not proved, for the price statistic.\hfill{\normalsize\itshape Thank you; questions are welcome.}}
\end{frame}

\end{document}
